\chapter{Queuing models}
\label{ch:colas}

\section{From flow to waiting}
Queuing theory relates variability, capacity, and congestion
\cite{gross2008fundamentals}. In Kendall notation, M denotes Poisson arrivals
or exponential times, depending on its position. Service rate \(\mu\) is
per server; mean service time is \(1/\mu\). All rates share a time unit.
QSB implements M/M/1 and a finite queue with multiple servers using a
birth-and-death model.

\section{M/M/1 in steady state}
The assumptions are independent arrivals, exponential service, one server,
and unlimited population and waiting space. Stability requires
\(\lambda<\mu\). With \(\rho=\lambda/\mu\):
\[
 P_0=1-\rho,\quad L=\frac{\rho}{1-\rho},\quad
 L_q=\frac{\rho^2}{1-\rho},\quad
 W=\frac1{\mu-\lambda},\quad W_q=\frac{\lambda}{\mu(\mu-\lambda)}.
\]
\(L\) includes customers being served; \(L_q\) includes only those waiting.
Little's law gives \(L=\lambda W\) and \(L_q=\lambda W_q\).
The backend reports \texttt{mm1ClosedForm}/\texttt{closedForm}. It does not
calculate the distribution or percentiles of an individual waiting time.

\subsection{Example: a service counter}
\path{manual/examples/prob_cola.json} specifies four arrivals/hour and six
services/hour. Run from the root:
\begin{lstlisting}
.build/debug/qsb validate-queuing-json manual/examples/prob_cola.json
.build/debug/qsb solve-queuing-json manual/examples/prob_cola.json
\end{lstlisting}
Utilization is \(2/3\), \(P_0=1/3\), \(L=2\), \(L_q=4/3\),
\(W=0.5\) hours, and \(W_q=1/3\) hours: thirty minutes in the system,
twenty waiting, and ten receiving service. Being busy two thirds of the time
does not mean that two thirds of customers wait exactly twenty minutes.
If arrivals reach six/hour, the model no longer has these steady-state means,
and QSB returns a stability diagnostic.

\section{Multiple servers and finite capacity}
With \(c\) servers and \(b\) waiting places, system capacity is \(K=c+b\).
State \(n\) has arrival rate \(\lambda\) when \(n<K\), and zero when
full; its departure rate is \(\min(n,c)\mu\).
The probabilities are normalized after being obtained from
\[
 p_n=p_0\prod_{i=1}^n\frac{\lambda}{\min(i,c)\mu}.
\]
Blocking probability is \(p_K\), the effective arrival rate
is \(\lambda_e=\lambda(1-p_K)\), and Little's law uses \(\lambda_e\),
rather than the offered rate. Because it is finite, this model can have a
stationary distribution even when offered demand exceeds its service capacity;
this produces blocking, rather than extra capacity.

QSB requires exponential arrivals in batches of one and an infinite customer population.
It uses mean service time to obtain the rate per server; nonexponential labels
produce warnings because their variability is not modeled. The document always
retains \texttt{approximate} for this route. A case with normal service times
must not be presented as an exact M/G/c/K solution.

\subsection{Example: two counters and three waiting places}
\path{prob_cola_finita.json} retains four arrivals/hour and a mean service
time of ten minutes, with two servers and three waiting places. Validate and
solve using the same commands with the file changed. The results are:
\begin{center}
\begin{tabular}{lr}
\toprule
Metric & Value \\
\midrule
Total capacity & 5 \\
Blocking probability & 0.004124 \\
Effective arrivals per hour & 3.983505 \\
Average utilization per server & 0.331959 \\
Mean number in the system & 0.738144 \\
Mean number waiting & 0.074227 \\
Time in the system (hours) & 0.185300 \\
Waiting time (hours) & 0.018634 \\
\bottomrule
\end{tabular}
\end{center}
Mean waiting time is approximately \(1.118012\) minutes among admitted
customers. The total utilization of both servers corresponds to \(0.663918\)
busy servers, rather than the individual utilization percentage.

\section{Costs, interface, and limitations}
Optional costs distinguish busy and idle servers, waiting customers, and
customers in service; the finite queue adds rejection and waiting places.
State costs are weighted by expected occupancy; rejection cost is multiplied
by the rate of blocked customers. The total is a cost per time unit, rather
than per customer.

In macOS, the creator offers both variants. Enter rates for M/M/1, or mean
times, servers, and waiting space for a finite queue; validate and solve. The
metrics panel shows utilization, effective flow, blocking, numbers, and times.
The finite variant also shows state probabilities, highlighting the full state;
optional costs are itemized. Check the time unit and assumptions alongside the
results. Normalized JSON is the alternative for inspection and export.

Use \path{solve-queuing-json} to accept normalized input.

The historical commands \path{solve-mm1-json} and
\path{solve-finite-queue-json} accept a legacy file despite their suffix:
the suffix indicates JSON output. Always use the generic route to exchange
normalized inputs. Native and validation backends are available; no external
solver is installed for this family.

Priorities, customer abandonment due to impatience, finite customer populations,
general batch arrivals, and rate changes by time of day are not modeled.
The simulation in the Markov and simulation chapter allows a different class
of assumptions to be explored, but its estimates do not replace these
steady-state formulas.

\section{Exercises}
\begin{enumerate}
\item Change \(\lambda\) to 5.5 in M/M/1 and compare utilization and waiting.
Explain why the growth in waiting time is nonlinear.
\item Remove the waiting places from the finite queue. Check that blocking
occurs without a positive queue length, and use \(\lambda_e\) to verify
Little's law.
\item Compare adding a server with adding five waiting places: which changes
service capacity, and which changes the space for customers?
\end{enumerate}
